\documentclass[11pt]{article}
\usepackage[margin=1in]{geometry}
\usepackage{amsmath,amssymb,amsthm,mathtools}
\usepackage[T1]{fontenc}
\usepackage{lmodern}
\usepackage{microtype}
\usepackage[colorlinks=true,urlcolor=blue,citecolor=blue,linkcolor=blue]{hyperref}
\setlength{\parskip}{3pt}
\setlength{\parindent}{1.2em}
\newtheorem{theorem}{Theorem}
\newtheorem{lemma}{Lemma}
\newtheorem{proposition}{Proposition}
\theoremstyle{remark}\newtheorem*{remark}{Remark}
\newcommand{\e}{\mathrm e}
\newcommand{\T}{\mathbb T}
\newcommand{\eps}{\varepsilon}
\newcommand{\M}{\mathfrak M}
\newcommand{\m}{\mathfrak m}
\newcommand{\SSS}{\mathfrak S}
\DeclareMathOperator{\meas}{meas}
\title{A short modern deduction of \(G(4)=16\)\\
\large Translation, sampling, and a smooth ordinary cutoff}
\author{Research draft}
\date{22 November 2026}
\begin{document}
\maketitle
\begin{abstract}
We deduce Davenport's theorem \(G(4)=16\) from the quartic Vinogradov mean
value theorem and Hua's eighth-moment estimate. We prove the needed
minor-arc transfer by translating the summation interval and sampling a
common dyadic majorant. A fixed smooth cutoff then permits a direct
Poisson-summation treatment of the same major arcs. We include the uniform
local-density argument, especially at the prime~2. No new Waring bound or
mean-value theorem is claimed.
\end{abstract}

\section{Statement and established inputs}

Let \(G(k)\) be the least integer \(s\) for which every sufficiently large
positive integer is a sum of at most \(s\) nonnegative integral \(k\)th powers.
We prove the following classical result of Davenport \cite{Davenport}.
\begin{theorem}\label{thm:main}
One has \(G(4)=16\).
\end{theorem}

The proof uses two established mean-value estimates. Write
\(\e(t)=\exp(2\pi it)\) and \(\T=\mathbb R/\mathbb Z\). The quartic case of
the main conjecture in Vinogradov's mean value theorem
\cite[Theorem~1.1]{BDG} gives
\begin{equation}\label{eq:vmvt}
 \int_{\T^4}\left|\sum_{1\le x\le X}
 \e(\beta_1x+\beta_2x^2+\beta_3x^3+\beta_4x^4)\right|^{20}
 d\boldsymbol\beta\ll_\eps X^{10+\eps}.
\end{equation}
Hua's lemma \cite[Lemma~2.5]{Vaughan} gives
\begin{equation}\label{eq:hua}
 \int_\T\left|\sum_{1\le x\le X}\e(\alpha x^4)\right|^8d\alpha
 \ll_\eps X^{5+\eps}.
\end{equation}
The estimate \eqref{eq:vmvt} holds for an interval of \(X\) consecutive
integers as well: changing the origin of \(x\) induces an invertible
integral triangular change of the four coefficient variables on \(\T^4\).

The transfer mechanism is the translation principle behind
Wooley's minor-arc mean-value theorem \cite[Theorem~2.1]{Wooley}; its
published proof has a corrigendum \cite{Wooleycorr}. We give the required
specialization directly below. Dyadic maximal majorants, sampling, and
translation orbits also appear in the recent preprint \cite[Sections~3--5]{Hu};
no result from that preprint is used here. The present goal is a compact
deduction with all remaining hypotheses visible, rather than a claim of
priority for these devices.

Fix a nonnegative \(w\in C_c^\infty((0,1))\), supported in
\([3/8,5/8]\), and positive near \(1/2\). For example, take
\[
 w(t)=
 \begin{cases}
 \exp\!\left(-\dfrac1{1-64(t-1/2)^2}\right),&|t-1/2|<1/8,\\
 0,&|t-1/2|\ge1/8.
 \end{cases}
\]
For a positive integer \(n\), put \(P=n^{1/4}\) and
\[
 f(\alpha)=\sum_{x\in\mathbb Z}w(x/P)\e(\alpha x^4),\qquad
 R_w(n)=\int_\T f(\alpha)^{16}\e(-n\alpha)d\alpha.
\]
Orthogonality expresses \(R_w(n)\) as a sum of nonnegative weights over
representations of \(n\) by sixteen positive fourth powers. Hence
\(R_w(n)>0\) suffices for the upper bound in Theorem~\ref{thm:main}.
Constants below may depend on \(w\) and \(\eps>0\), which are fixed.

\section{An elementary sampling transfer}

We first record the sampling estimate used in the argument.
\begin{lemma}\label{lem:sampling}
Suppose that \(T\) is a nonnegative trigonometric polynomial of degree at
most \(D\ge1\), and that \(t_1,\ldots,t_Y\in\T\) are separated in circular
distance by at least \(\delta\). Then
\[
 \sum_{y=1}^Y T(t_y)\ll (D+\delta^{-1})\int_\T T(t)dt.
\]
\end{lemma}
\begin{proof}
Take disjoint intervals \(I_y\) of length \(\delta\), centered at \(t_y\).
For \(t\in I_y\), the fundamental theorem of calculus gives
\(T(t_y)\le T(t)+\int_{I_y}|T'|\). Integrating and summing yields
\[
 \sum_yT(t_y)\le\delta^{-1}\int_\T T+\int_\T|T'|.
\]
The elementary bound \(\|T'\|_1\ll D\|T\|_1\) finishes the proof.
For completeness, let \(K_N(t)=N^{-1}(\sin(\pi Nt)/\sin(\pi t))^2\)
be the Fej\'er kernel. The kernel \(U_D=2K_{2D}-K_D\) reproduces every
polynomial of degree at most \(D\). Direct differentiation gives
\(\|K_N'\|_1\ll N\): use \(|K_N'|\ll N^2\) for \(\|t\|\le N^{-1}\),
and \(|K_N'|\ll\|t\|^{-2}+N^{-1}\|t\|^{-3}\) elsewhere.
Thus \(T'=T*U_D'\) has the asserted \(L^1\) bound.
\end{proof}

For \(P^{1/2}\le Q\le P\), define the disjoint major arcs
\begin{equation}\label{eq:arcs}
 \M(Q)=\bigcup_{\substack{q\le Q\\(a,q)=1}}
 \{\alpha\in\T:|q\alpha-a|\le QP^{-4}\},\qquad
 \m(Q)=\T\setminus\M(Q).
\end{equation}
Distances in this notation are understood on the circle. Disjointness
follows from \(2Q^2<P^4\) for large \(P\).

\begin{proposition}\label{prop:transfer}
For every \(\eps>0\),
\[
 \int_{\m(Q)}|f(\alpha)|^{20}d\alpha
 \ll_\eps P^{17+\eps}Q^{-2}.
\]
\end{proposition}
\begin{proof}
Put
\[
 F(\beta_1,\beta_2,\alpha)
 =\sum_xw(x/P)\e(\beta_1x+\beta_2x^2+\alpha x^4),
 \qquad I(\alpha)=\int_{\T^2}|F|^{20}d\beta_1d\beta_2.
\]
The polynomial \(F^{10}\) has at most \(10P+1\) possible frequencies in
\(\beta_1\), and at most \(10P^2+1\) in \(\beta_2\). Applying
Cauchy--Schwarz to the Fourier coefficients in each variable gives
\begin{equation}\label{eq:evaluate}
 |f(\alpha)|^{20}\ll P^3 I(\alpha).
\end{equation}

Choose a power of two \(N\) with \(2P\le N<4P\), and let \(\mathcal D\)
be all dyadic subintervals of \(\{1,\ldots,N\}\). Put
\[
 T_J(\boldsymbol\beta)=\sum_{x\in J}
 \e(\beta_1x+\beta_2x^2+\beta_3x^3+\beta_4x^4),\qquad
 H(\boldsymbol\beta)=\sum_{J\in\mathcal D}|T_J(\boldsymbol\beta)|^{20}.
\]
Every interval is a disjoint union of at most \(2(1+\log_2N)\) dyadic
intervals. Consequently
\begin{equation}\label{eq:dyadic}
 \max_{J\text{ an interval in }[1,N]}|T_J(\boldsymbol\beta)|^{20}
 \ll(\log P)^{19}H(\boldsymbol\beta).
\end{equation}
By \eqref{eq:vmvt} and its translation invariance,
\begin{equation}\label{eq:hmean}
 \int_{\T^4}H\ll_\eps
 \sum_{j=0}^{\log_2N}\frac{N}{2^j}(2^j)^{10+\eps}
 \ll_\eps P^{10+\eps}.
\end{equation}

Let \(Y=\lfloor Q/8\rfloor\). For \(\alpha\in\m(Q)\), the points
\(-4y\alpha\pmod1\), \(1\le y\le Y\), are \(\delta\)-separated with
\(\delta=QP^{-4}\). Indeed, otherwise \(h=|y-z|\) satisfies
\(\|4h\alpha\|\le QP^{-4}\); reducing the nearest rational with
denominator \(4h\le Q\) contradicts \(\alpha\notin\M(Q)\).

In \(F\), substitute \(x=z-y\). Apart from a constant phase, the new
coefficients are
\[
 \beta_1'=\beta_1-2y\beta_2-4y^3\alpha,
 \quad\beta_2'=\beta_2+6y^2\alpha,
 \quad\beta_3'=-4y\alpha,
 \quad\beta_4'=\alpha.
\]
The first two coefficients are related to the old ones by a
measure-preserving integral shear and a translation on \(\T^2\).
The translated weight \(w((z-y)/P)\) has bounded total variation,
uniformly in \(y\). Abel summation bounds its weighted exponential sum
by \(C_w\) times the largest interval sum in \eqref{eq:dyadic}.
Writing
\[
 H_\alpha(t)=\int_{\T^2}H(\beta_1,\beta_2,t,\alpha)d\beta_1d\beta_2,
\]
we obtain, for every \(1\le y\le Y\),
\(I(\alpha)\ll_w(\log P)^{19}H_\alpha(-4y\alpha)\).

Now \(H_\alpha\) is a nonnegative trigonometric polynomial in \(t\)
of degree at most \(10N^3\ll P^3\). Average over \(y\), and apply
Lemma~\ref{lem:sampling}. Since \(Y\asymp Q\) and \(Q\le P\),
\[
 I(\alpha)\ll_w
 \frac{(\log P)^{19}}Q(P^3+P^4/Q)\int_\T H_\alpha(t)dt
 \ll_w\frac{P^4(\log P)^{19}}{Q^2}\int_\T H_\alpha(t)dt.
\]
Combine this with \eqref{eq:evaluate}, integrate in \(\alpha\), enlarge
the last domain to \(\T^4\) by nonnegativity, and use
\eqref{eq:hmean}. Absorbing the logarithm into \(P^\eps\) proves the
proposition.
\end{proof}

Expanding the weighted eighth moment and using orthogonality shows that
it is at most \(\|w\|_\infty^8\) times the unweighted moment in
\eqref{eq:hua}. H\"older's inequality therefore gives
\begin{equation}\label{eq:minor16}
 \int_{\m(Q)}|f|^{16}
 \le\left(\int_{\m(Q)}|f|^{20}\right)^{2/3}
     \left(\int_\T|f|^8\right)^{1/3}
 \ll_\eps P^{13+\eps}Q^{-4/3}.
\end{equation}
From now on take
\begin{equation}\label{eq:qchoice}
 Q=P^{7/8}.
\end{equation}
The right side of \eqref{eq:minor16} is \(P^{71/6+\eps}=o(P^{12})\).

\section{The major arcs by Poisson summation}

Define
\[
 S(q,a)=\sum_{r\pmod q}\e(ar^4/q),\qquad
 v(\beta)=\int_{\mathbb R}w(x/P)\e(\beta x^4)dx.
\]
We shall use
\begin{equation}\label{eq:completebound}
 |S(q,a)|\ll_\eps q^{3/4+\eps}\quad((a,q)=1),
\end{equation}
proved in Section~\ref{sec:complete}. Since the support of \(w\) is
bounded away from zero, repeated integration by parts gives
\begin{equation}\label{eq:vbound}
 |v(\beta)|\ll_{A,w}P(1+P^4|\beta|)^{-A}
 \quad\text{for every fixed }A>0.
\end{equation}

For \(\alpha=a/q+\beta\in\M(Q)\), Poisson summation in residue classes
modulo \(q\) yields
\[
 f(\alpha)=\frac1q\sum_{m\in\mathbb Z}S(q,a,m)
 P\int_{\mathbb R}w(t)\e(\beta P^4t^4-mPt/q)dt,
\]
where \(S(q,a,m)=\sum_{r\pmod q}\e((ar^4+mr)/q)\), up to a harmless
sign change of \(m\). The zero mode is \(q^{-1}S(q,a)v(\beta)\).
For \(m\ne0\), the first derivative of the phase in \(t\) has modulus
at least \(|m|P/(2q)\) for large \(P\): its quartic contribution is
\(O(Q/q)\), whereas its linear contribution has size \(|m|P/q\).
The normalized higher derivatives are bounded uniformly. Repeated
integration by parts and the trivial bound \(|S(q,a,m)|\le q\) give
\begin{equation}\label{eq:poisson}
 f(\alpha)=V(\alpha)+O_{A,w}\bigl(P(Q/P)^A\bigr),
 \qquad V(\alpha)=q^{-1}S(q,a)v(\beta).
\end{equation}
Take \(A=32\). By \eqref{eq:qchoice}, the error is \(O_w(P^{-3})\).
Also \(|f|+|V|\ll_wP\) and
\(\meas\M(Q)\ll Q^2P^{-4}\). Thus replacing \(f^{16}\) by
\(V^{16}\) in the major-arc integral has total error \(O_w(P^{10})\).

Put
\[
 A(q,n)=q^{-16}\sum_{\substack{a\pmod q\\(a,q)=1}}
 S(q,a)^{16}\e(-an/q),\qquad
 \SSS(n)=\sum_{q\ge1}A(q,n),
 \quad J(n)=\int_{\mathbb R}v(\beta)^{16}\e(-n\beta)d\beta.
\]
The bound \eqref{eq:completebound} gives
\(|A(q,n)|\ll_\eps q^{-3+\eps}\), uniformly in \(n\). In particular,
the singular series is absolutely convergent. Equations
\eqref{eq:vbound} and \eqref{eq:completebound} give
\begin{equation}\label{eq:majmain}
 \int_{\M(Q)}V(\alpha)^{16}\e(-n\alpha)d\alpha
 =\SSS(n)J(n)+O_\eps(P^{12}Q^{-2+\eps}).
\end{equation}
For example, the weaker bound
\(|v(\beta)|\ll P(1+P^4|\beta|)^{-1/4}\) already suffices:
extending the interval for denominator \(q\le Q\) costs
\(O(P^{12}(Q/q)^{-3})\). Summing against \(|A(q,n)|\), and then
adding the tail \(q>Q\), proves \eqref{eq:majmain}.

Let \(h(u)=\tfrac14u^{-3/4}w(u^{1/4})\) for \(u>0\), and zero
otherwise. Scaling and Fourier inversion give
\[
 J(n)=C_wP^{12},\qquad C_w=(h*\cdots*h)(1)>0,
\]
with sixteen factors. Positivity holds because \(h\) is positive near
\(1/16\): choose its first fifteen arguments sufficiently near
\(1/16\), and the final argument determined by their sum is still in
that positive neighborhood. Combining \eqref{eq:minor16},
\eqref{eq:poisson}, and \eqref{eq:majmain}, we obtain
\begin{equation}\label{eq:weightedcount}
 R_w(n)=C_w\SSS(n)n^3+O_{w,\eps}(n^{71/24+\eps}).
\end{equation}
It remains to establish a lower bound for \(\SSS(n)\) uniform in \(n\).

\section{Complete sums and uniform local positivity}\label{sec:complete}

For an odd prime \(p\), the multiplicative-character expression for the
quartic sum contains at most three nontrivial Gauss sums; each has
modulus \(\sqrt p\), by character orthogonality. Thus
\begin{equation}\label{eq:gauss}
 |S(p,a)|\le3\sqrt p\qquad(p\nmid a).
\end{equation}
For \(h\ge2\), summing first over shifts by \(p^{h-1}\) cancels
every term with \(p\nmid x\), since its linear coefficient is
\(4ax^3\not\equiv0\pmod p\). The remaining multiples of \(p\) give
\begin{equation}\label{eq:oddrec}
 S(p^h,a)=p^{h-1}\ (2\le h\le4),\qquad
 S(p^h,a)=p^3S(p^{h-4},a)\ (h\ge5).
\end{equation}
Together with \eqref{eq:gauss}, this gives
\(|S(p^h,a)|\le3p^{3h/4}\).

For \(p=2\) and \(h\ge5\), sum in cosets generated by the shift
\(2^{h-3}\). At each step the fourth-power phase changes by
\(2^{h-1}x^3\) modulo \(2^h\), so odd terms cancel. Therefore
\[
 S(2^h,a)=8S(2^{h-4},a)\quad(h\ge5).
\]
For \(1\le h\le4\), separation into odd and even terms gives
\(S(2^h,a)=2^{h-1}(1+\e(a/2^h))\), whence
\(|S(2^h,a)|\le2\cdot2^{3h/4}\). Chinese-remainder factorization
now proves \eqref{eq:completebound}, since
\(3^{\omega(q)}\ll_\eps q^\eps\).

Let \(M_h(n;p)\) count the vectors modulo \(p^h\) satisfying
\(x_1^4+\cdots+x_{16}^4\equiv n\pmod{p^h}\). Orthogonality,
grouped according to the conductor of the additive character, gives
\begin{equation}\label{eq:localcount}
 p^{-15h}M_h(n;p)=\sum_{j=0}^h A(p^j,n).
\end{equation}
The Chinese remainder theorem makes \(A(q,n)\) multiplicative in
\(q\). Absolute convergence therefore gives
\begin{equation}\label{eq:localprod}
 \SSS(n)=\prod_p\sigma_p(n),\qquad
 \sigma_p(n)=\lim_{h\to\infty}p^{-15h}M_h(n;p).
\end{equation}

For \(p=3,5,7,11\), represent any residue modulo \(p\) by its
representative in \(\{1,\ldots,p\}\) ones, padding with zeros to
sixteen terms. This provides a nonzero coordinate. For \(p\ge13\),
\eqref{eq:gauss} gives
\[
 M_1(n;p)\ge p^{15}-3^{16}p^8>1,
 \qquad 13^7>3^{16}.
\]
There is again a nonzero coordinate, even when \(n\equiv0\pmod p\).
The derivative \(4x_i^3\) at that coordinate is a unit. Ordinary
Hensel lifting supplies \(p^{15}\) lifts at each step, so
\begin{equation}\label{eq:oddlocal}
 \sigma_p(n)\ge p^{-15}>0\qquad(p\text{ odd}).
\end{equation}

At the prime~2, every integer congruent to~1 modulo~16 has an odd
fourth root modulo \(2^h\) for every \(h\ge4\). An elementary
induction uses
\[
 (x+2^{h-2})^4\equiv x^4+2^hx^3\pmod{2^{h+1}}
 \qquad(h\ge4,\ x\text{ odd}),
\]
which toggles the next bit. Choose the representative \(r\) of
\(n\pmod{16}\) in \(\{1,\ldots,16\}\). Require one distinguished
coordinate only to be odd; prescribe \(r-1\) of the other fifteen
coordinates to be~1 modulo~16 and the remaining \(16-r\) to be~0.
Every choice of these fifteen variables with their prescribed residues
can be completed, because the needed
fourth power is~1 modulo~16. Hence
\begin{equation}\label{eq:twolocal}
 M_h(n;2)\ge2^{15(h-4)},\qquad \sigma_2(n)\ge2^{-60}.
\end{equation}
This includes every \(n\) divisible by~16, with no descent issue.

Finally, \eqref{eq:gauss} bounds the \(h=1\) term of an odd local
factor by \(O(p^{-7})\), and \eqref{eq:oddrec} bounds its terms
\(h\ge2\) by \(O(p^{-3h})\). Uniformly in \(n\),
\[
 \sigma_p(n)=1+O(p^{-6}).
\]
The product over all sufficiently large primes is therefore bounded
below by a fixed positive constant. Equations \eqref{eq:oddlocal}
and \eqref{eq:twolocal} handle the remaining finite set. Thus
\(\inf_{n\ge1}\SSS(n)>0\). The asymptotic formula
\eqref{eq:weightedcount} proves \(G(4)\le16\).

\section{The obstruction to fifteen variables}

For every integer \(a\ge0\), the number \(31\cdot16^a\) is not a
sum of at most fifteen nonnegative fourth powers. If \(a\ge1\),
reduction modulo~16 forces every variable to be even: a fourth power
is~0 or~1 modulo~16, and there are fewer than sixteen variables.
Divide by~16 and repeat, reaching~31. Its possible fourth-power
summands are only~0,~1, and~16. Without a~16 the sum is at most~15;
with one~16 it needs fifteen more ones, hence sixteen terms; two~16s
are too large. This proves \(G(4)\ge16\) and completes the theorem.

\begin{remark}
The argument proves a known conclusion from deep established input;
it is not an elementary proof from first principles. It uses no
hypothesis concerning zeros of zeta or \(L\)-functions. An alternative
unsmoothed version gives
\[
 R_{16,4}(n)=\frac{\Gamma(5/4)^{16}}{\Gamma(4)}\SSS(n)n^3
 +O_\eps(n^{83/28+\eps})
\]
by choosing \(Q=P^{6/7}\) and balancing the elementary Euler-summation
error. The weighted version above is shorter when existence is the
only target. Bourgain \cite[Section~2, Theorem~3]{Bourgain} already
deduced the stronger known fifteen-variable asymptotic formula from
modern VMVT and Wooley's framework; that result is not being claimed
anew here.
\end{remark}

\begin{thebibliography}{9}
\bibitem{Davenport}
H. Davenport, \emph{On Waring's problem for fourth powers},
Ann. of Math. (2) \textbf{40} (1939), 731--747.

\bibitem{BDG}
J. Bourgain, C. Demeter, and L. Guth,
\emph{Proof of the main conjecture in Vinogradov's mean value theorem
for degrees higher than three}, Ann. of Math. (2) \textbf{184} (2016),
633--682. \url{https://arxiv.org/abs/1512.01565}.

\bibitem{Vaughan}
R. C. Vaughan, \emph{The Hardy--Littlewood Method}, second edition,
Cambridge University Press, 1997, Lemma~2.5.

\bibitem{Wooley}
T. D. Wooley, \emph{The asymptotic formula in Waring's problem},
Int. Math. Res. Not. IMRN (2012), no.~7, 1485--1504.
\url{https://www.math.purdue.edu/~twooley/publ/2012asympwp.pdf}.

\bibitem{Wooleycorr}
T. D. Wooley, \emph{Corrigendum: The asymptotic formula in Waring's
problem}, Int. Math. Res. Not. IMRN (2015), no.~20, 10702.
\url{https://www.math.purdue.edu/~twooley/publ/2012asympwpcorr.pdf}.

\bibitem{Bourgain}
J. Bourgain, \emph{On the Vinogradov mean value}, 2016,
Section~2, Theorem~3. \url{https://arxiv.org/abs/1601.08173}.

\bibitem{Hu}
X. Hu, \emph{Improved Weyl bounds on short intervals},
arXiv preprint, 2026, Sections~3--5.
\url{https://arxiv.org/abs/2609.02478}.
\end{thebibliography}
\end{document}
